% RecSys Odyssey repository-backed lecture note.
% Add figures with: \coursefigure{figures/name.svg}{Caption}
\section{Protect the whole ecosystem}

\begin{abstract}
Healthy systems create value for users without silently starving parts of the catalogue.
\end{abstract}

\subsection{Concept}
Ecosystem monitoring segments outcomes by user group, item age, creator size and catalogue region. It tracks exposure concentration, supplier coverage, satisfaction and safety beside core engagement. Constraints can reserve exploration, cap repetition and protect eligibility, while human review handles value judgments metrics cannot settle. Fairness is not one universal number: teams must state who is affected, what allocation is being compared and which harms are unacceptable.

\begin{equation*}
launch only if primary value ↑ and user, safety, fairness, catalogue guardrails hold
\end{equation*}

\subsection{Key terms}
\begin{itemize}
\item \textbf{Exposure parity.} A declared comparison of how attention is allocated across groups.
\item \textbf{Catalogue health.} The breadth, freshness and sustainable supply of recommendable items.
\item \textbf{Guardrail.} A hard or monitored boundary protecting against unwanted side effects.
\end{itemize}
